\section*{Guide to the strengthened proof}
\addcontentsline{toc}{section}{Guide to the strengthened proof}
\label{sec:proof-architecture-near-root-v4}

The argument has three stages---location, existence, and amplification. The
new simplification occurs at their interface: the signed profile needs only a
one-part buffer above its continuous root. This retains the entire leading
root separation while leaving every second-moment estimate valid.

\subsection*{1. Root separation and minimal integer buffering}

Let $r_+(n)$ be the continuous root of the unrestricted coloring objective and
$r_4^{\mathrm{co}}(n)$ the root of the signed objective on deficits
$\{2,3,4,5\}$. The phase and finite-dual packages give, uniformly in the
complete phase,
\[
  r_+-r_4^{\mathrm{co}}
  =
  \left[
    \frac{(\log 2)^2}{4}A_4(\delta_n)+o(1)
  \right]
  \frac{n}{(\log n)^3},
  \qquad
  A_4(\delta)=\log 2-D_4(\delta).
\]
Instead of taking the midpoint, define
\[
  k_{\mathrm{co}}=\lceil r_4^{\mathrm{co}}\rceil+1.
\]
The displacement $d_n=k_{\mathrm{co}}-r_4^{\mathrm{co}}$ lies in $[1,2)$.
Moving by fewer than two classes changes $n/k$ by
$O((\log n)^2/n)$, so the full segment from the real root to
$k_{\mathrm{co}}$ stays in the same derivative corridor. The slope theorem
therefore gives
\[
  \Phi_n(k_{\mathrm{co}})
  =\Theta((\log n)^2),
  \qquad
  \Phi_n(k)=L_{S_4}(n,k)+(\log 2)k.
\]
After exact tangent rounding and Stirling extraction,
\[
  \log\mathbb E Z_{\mathbf k}^{\mathrm{sgn}}
  =\Theta((\log n)^2).
\]
The profile remains uniformly positive in all four coordinates and satisfies
both integer conservation laws exactly.

Because $k_{\mathrm{co}}=r_4^{\mathrm{co}}+O(1)$ and the chromatic threshold is
$r_++O(\log n)$, the deterministic location gap retains the full coefficient:
\[
  k_\chi^- - k_{\mathrm{co}}
  =
  \left[
    \frac{(\log 2)^2}{4}A_4(\delta_n)+o(1)
  \right]
  \frac{n}{(\log n)^3}.
\]

\subsection*{2. Why the smaller first moment is sufficient}

The partial-diagonal proof is the only part of the normalized second moment
that uses the size of the signed first moment.

In the central range, equation (7.17) divides its logarithm by
$k_{\mathrm{co}}=\Theta(n/\log n)$. Thus
\[
  \frac{\log\mathbb E Z}{k_{\mathrm{co}}}
  =O\!\left(\frac{(\log n)^3}{n}\right)=o(1),
\]
which is exactly what the affine extraction requires.

In the full corner, the reverse recurrence gives $B(h)\le1$, hence
\[
  D(k-h)
  \le
  \bigl(\mathbb E Z_{\mathbf k}^{\mathrm{sgn}}\bigr)^{-1}
  \le e^{-c(\log n)^2}.
\]
There are at most $(k_{\mathrm{co}}+1)^4$ residual profiles, whose logarithm
is $O(\log n)$. Therefore their total contribution tends to zero. The
midpoint-scale margin $\Theta(n/\log n)$ was stronger than needed by a wide
margin.

The empty corner, central rate negativity, endpoint transport, and residual
attachment use only the same phase corridor, coordinate positivity, and
$k_i=\Theta(n/\log n)$. They are unchanged.

\subsection*{3. Fusing deficits with endpoint transport}

For endpoint type $(i,j)$, let
\[
  B_{ij}=(1-\rho_{ij})^{-1},
  \qquad
  \widetilde Q_{ij}=B_{ij}Q_{ij}.
\]
The complete positive-deficit fiber is summed geometrically before endpoint
tables are aggregated. Multiplying the square-free transport inequality by the
table factor $(B^L)^2$ and applying arithmetic--geometric mean gives
\[
  \operatorname{BareSkeletonSum}_n
  \le
  \Sigma_n^K\sum_rD(r),
  \qquad
  \Sigma_n:=\max_i\sum_j\widetilde Q_{ij}.
\]
The exact phase identity yields
\[
  \rho_{16}=O\!\left(\frac{(\log n)^{5/2}}{\sqrt n}\right),
  \qquad
  \Sigma_n=1+O\!\left(\frac{(\log n)^{5/2}}{\sqrt n}\right),
\]
and hence
\[
  \log\operatorname{BareSkeletonSum}_n
  =O\!\left(\sqrt n(\log n)^{3/2}\right).
\]
One kernel theorem now consumes the deficit product, support size, endpoint
regrouping, and transport row sum.

\subsection*{4. Critical-quarter residual attachments}

For residual mass $m_0$ and class-size cap $U$, use the threshold
\[
  T_U=64U2^{U/4}.
\]
When $m_0\ge T_U$, every residual intensity satisfies
\[
  \theta_{ab}
  \le\frac{\mathrm eU}{64}2^{-U/4}.
\]
The log-convex endpoint estimate then gives
$q_{ab}\le C\theta_{ab}^2$, and the exact square-intensity identity yields an
attachment factor $e^{O(U^2)}$. When $m_0<T_U$, the crude factor
$2^{Um_0/2}$ gives
\[
  \log\mathcal A(M,j)
  =O(U^2 2^{U/4}).
\]
Since $2^U=\Theta(n^2/(\log n)^2)$, both the fused skeleton and residual
attachment costs are bounded by
\[
  O\!\left(\sqrt n(\log n)^{3/2}\right).
\]
Thus
\[
  1
  \le
  \frac{\mathbb E Z^2}{(\mathbb E Z)^2}
  \le
  \exp\!\left(C_\sharp\sqrt n(\log n)^{3/2}\right).
\]

\subsection*{5. Explicit amplification loss}

Paley--Zygmund supplies a seed with probability at least
\[
  \exp\!\left\{-C_\sharp\sqrt n(\log n)^{3/2}\right\}.
\]
Applying the arbitrary-seed amplifier with $r=\log n$ adds
\[
  O\!\left(n^{3/4}(\log n)^{-1/4}\right)
  =o\!\left(\frac{n}{(\log n)^3}\right)
\]
classes with failure probability tending to zero. The amplifier therefore
consumes none of the leading root-gap coefficient.

Finally, the exact four-support certificate
\[
  A_4(\delta)
  >
  \log\!\left(\frac{20000}{12777}\right)
\]
contains fixed slack over $\log(1000/639)$. The retained uniform coefficient is
\[
  \frac{(\log 2)^2}{4}
  \log\!\left(\frac{1000}{639}\right)
  =0.053792819616758\ldots.
\]
